\subsection{Significance of the Project}

\subsubsection{Scientific Perspectives}
From a scientific standpoint, aBridgeAI makes several contributions to the intersection of educational technology, cognitive science, and applied machine learning.

The primary scientific contribution is the formalisation of an algorithmically derived quality response value for the SM-2 spaced repetition algorithm in institutional deployment contexts. Standard SM-2 relies on learner self-assessment as its sole input signal, an assumption that holds in solo study environments but breaks down when the quality score influences externally visible progression gates. This work derives $Q$ from three objective, system-observable signals — response correctness, hint usage, and a dimensionless time ratio $\rho = T_{\text{actual}} / T_{\text{exp}}$ — replacing a subjective input with a deterministic computation grounded in response latency theory \cite{benjamin1996}. The threshold boundaries in the $Q$ assignment rule are derived as natural multiples of the teacher-defined expected response time rather than as independently calibrated constants, eliminating arbitrary parameterisation from the scheduling model.

A secondary contribution is the introduction of a calibration weight $\alpha$ applied asymmetrically to early EF increments, addressing the known cold-start overestimation problem in SM-2 initialisation. By dampening positive EF drift during the first three repetitions while applying negative increments at full magnitude, the model converges onto a student's true retention profile faster than the standard initialisation assumption permits. This mechanism is formalised within the SM-2 framework without requiring the full parametric complexity of successor algorithms such as FSRS \cite{ye2022fsrs}.

At the assessment layer, the work contributes a teacher-defined outcome architecture for AI-conducted mock interviews that decouples approved question-bank selection, live progression, and terminal evaluation. Rather than scoring open-ended responses against a model answer key, the system uses teacher-specified outcomes, server-controlled state, and a durable transcript. The target native runtime may use an in-session LLM sufficiency probe for provisional progression, while asynchronous terminal evaluation independently determines the recorded outcome verdict. This separation preserves open-ended interaction without giving the conversational model sole authority over assessment.

Finally, the system introduces three formally defined learner-facing metrics — the Knowledge Retention Estimate $\hat{R}_{s,\ell}$, Progression Readiness, and Review Compliance Rate $\rho_{s,\ell}$ — that together characterise a student's cognitive standing within a lesson from complementary diagnostic, gate, and engagement perspectives. The formal definitions of these metrics, their derivation from per-student SM-2 state, and their relationship to the no-decay design decision constitute a contribution to the literature on competency-based progression systems in spaced repetition contexts.

\subsubsection{Practical Perspectives}
From a practical standpoint, aBridgeAI addresses a concrete and growing problem in institutional education: the inability of existing tools to bridge the gap between passive knowledge consumption and verified competency demonstration at scale.

For \textbf{teachers and instructors}, the platform eliminates the manual overhead of quiz creation by automating question generation from uploaded course materials, while preserving full editorial control through a mandatory review and configuration phase. Teachers can define expected response times and EF thresholds per card, configure lesson unlock criteria, and specify interview outcome requirements — all without requiring knowledge of the underlying scheduling algorithm. The monitoring dashboard surfaces per-student and cohort-level retention metrics, enabling data-informed intervention rather than reliance on periodic examination snapshots.

For \textbf{students}, the platform replaces fixed assessment schedules with individually paced review cycles governed by demonstrated memory stability. Students are notified when cards are due for review at algorithmically optimal intervals, reducing the cognitive burden of self-scheduling while reinforcing retention through spaced practice. The interview module provides repeated access to structured practice sessions in a low-stakes environment, building the practical communication and problem-solving skills that quiz-based assessment cannot evaluate. Career Readiness reporting gives students a transparent, continuously updated signal of their progression toward professional competency targets.

For \textbf{organisations and educational managers}, the system provides institution-level visibility into cohort progression through career-scoped reporting. Career Readiness scores aggregated across course completion percentages give managers a quantitative basis for identifying at-risk cohorts, evaluating curriculum effectiveness, and demonstrating learning outcomes to accreditation bodies or industry partners — without requiring manual data collection or post-hoc grade analysis.

From a broader industry perspective, aBridgeAI demonstrates that the cost and complexity barriers associated with AI-assisted personalised learning can be constrained architecturally: question banks are pre-generated and teacher-approved, real-time interview conversation is bounded by server-controlled state and tools, richer reconciliation and final evaluation are asynchronous, and scheduling is grounded in a deterministic algorithm. This separation limits live-model responsibility while retaining personalised interaction.